\section{Map from experimental claims to proofs and evidence}
\label{sec:seq-sim-proof-map}

Every mathematical claim in the experimental exposition belongs to a row
below. The numbered proofs are part of this standalone report; consulting
the companion analysis is optional. The primary references identify the
well-known results from which the proofs proceed. Definitions, model
choices and finite observed records have their own evidential status.

\begingroup\small
\setlength{\tabcolsep}{4pt}
\begin{longtable}{>{\raggedright\arraybackslash}p{.25\textwidth}
                 >{\raggedright\arraybackslash}p{.38\textwidth}
                 >{\raggedright\arraybackslash}p{.28\textwidth}}
\toprule
Exposition/claim & Proof or exact record classification & Established foundation\\
\midrule\endfirsthead
\toprule
Exposition/claim & Proof or exact record classification & Established foundation\\
\midrule\endhead
\bottomrule\endfoot
Abstract and Section~\ref{sec:seq-sim-questions}: standardization, fixed
starting values, three means and IQ units &
Assumptions/definitions are stated there. Their consequences are proved in
Propositions~\ref{seqp:recursion} and \ref{seqp:variance-decomposition};
the exact grid is in Section~\ref{seqp:exact-grid}. &
Moment identities and Gaussian transformations,
\ref{pf:moments}, \ref{pf:gaussian}
\citep{durrett2019,gallager2013}.\\\midrule
Section~\ref{sec:seq-sim-model}: factor construction, admissible finite
correlations and the fixed-correlation infinite domain &
Proposition~\ref{seqp:domain} proves the two eigenspaces, the PSD
inequalities and the nonnegative factor realization. &
Spectral theorem/positive square root, \ref{pf:spectral}
\citep{axler2024}; covariance expansion.\\\midrule
Section~\ref{sec:seq-sim-model}: posterior factor, whole future block,
independence and the probability-zero initial history &
Proposition~\ref{seqp:conditioning} proves the inverse, residual
independence, full conditional kernel, covariance matrix and continuous
mean-version uniqueness; it treats \(c=1\) directly on its support. &
Gaussian characteristic functions, regression residual and conditional
integral identity, \ref{pf:gaussian}, \ref{pf:conditioning}
\citep{gallager2013,durrett2019}.\\\midrule
Sections~\ref{sec:seq-sim-model}, \ref{sec:seq-sim-means}: every prediction,
inclusive centre, population centre, normality and SD zero &
Proposition~\ref{seqp:recursion} proves regression plus induction, the
inclusive numerator and denominator, Gaussian closure and constant-list
variance. The positive difference from the next prediction is computed
explicitly in Section~\ref{sec:seq-sim-means}. &
Conditional regression and moment identities,
\ref{pf:conditioning}, \ref{pf:moments}; elementary induction.\\\midrule
Sections~\ref{sec:seq-sim-model}, \ref{sec:seq-sim-sds}: actual population
and conditional mean SDs; dependence on total count &
Proposition~\ref{seqp:recursion} counts all diagonal and ordered
off-diagonal covariances; Proposition~\ref{seqp:variance-decomposition}
separates between-history and remaining variance. &
Covariance bilinearity, conditional expectation and total variance,
\ref{pf:moments}, \ref{pf:expectation}
\citep{durrett2019}.\\\midrule
Section~\ref{sec:seq-sim-monotonicity}: decreasing population and
increasing conditional SD at \(c=3/5\) &
The displayed finite-difference identity and its three strictly positive
coefficients prove the conditional assertion for all integer \(r,k\geq1\).
The population difference is negative from \(c+(1-c)/(r+k)\). &
Proved variance formulas; order-preservation of the nonnegative square root.
No inference from plotted ordering is required.\\\midrule
Section~\ref{sec:seq-sim-limits}: prediction and inclusive-mean limits,
both SD limits, limiting Gaussian experiments and common-factor uncertainty &
That section proves the exact \(1/N\) error, conditional variance
expansion, square-root inequality and the common-normal mean-square
coupling. Proposition~\ref{seqp:domain} establishes the permissible
fixed-correlation range. &
The proved finite targets, Gaussian affine closure and moment identities,
\ref{pf:gaussian}, \ref{pf:moments}.\\\midrule
Sections~\ref{sec:seq-sim-means}--\ref{sec:seq-sim-limits}: all rational
targets and displayed analytic decimal approximations &
Section~\ref{seqp:exact-grid} proves substitutions and supplies exact
positive-integer certificates for every table target rounding and the
six-decimal worked example. &
Affine variance scaling and the proved target formulas; rational
arithmetic and order of positive squares.\\\midrule
Sections~\ref{sec:seq-sim-design}--\ref{sec:seq-sim-sds}: sample SD,
theoretical MC SE, chi-square SD intervals and their scope &
\ref{sm:normal-sample} proves the normal/chi-square sampling laws;
\ref{sm:sd-interval} inverts the pivot. \ref{sm:multiple} proves the
multiple-comparison bounds without assuming independent case estimates. &
Orthogonal basis extension \citep{axler2024}, Gaussian block independence
\citep{gallager2013}, quantile inversion and the union bound.\\\midrule
Sections~\ref{sec:seq-sim-design}, \ref{sec:seq-sim-replication}: pooling,
four replicates, output counts and shared random numbers &
\ref{sm:pooling} proves count/sum/within-between-square identities;
Section~\ref{seqp:record-scope} proves the finite grid arithmetic. The
seed/chunk/sharing choices are implementation records. &
Finite-sum identities. Independent-observation sampling laws are stated
ideal assumptions, rather than consequences of seed values.\\\midrule
Tables, figures, generated discussion and retained Ada comparison &
Observed estimates, maxima, intervals, CSV coordinates, seeds and proof
counts are finite reproduction records. Their exact scope and the
counterexample to inferring universal correctness from tests appear in
Section~\ref{seqp:record-scope}. &
Documented Dask/SciPy interfaces \citep{daskArray,daskDelayed,scipystats};
source-pinned program evidence has its stated separate trust boundaries.\\\midrule
Section~\ref{sec:seq-sim-interpretation}: what the experiment establishes,
model adequacy and distinction between predictions and outcomes &
Propositions~\ref{seqp:recursion} and \ref{seqp:variance-decomposition}
prove the model conclusions. Section~\ref{seqp:record-scope} separates
finite observations from universal claims and empirical validation. &
The stated Gaussian model and proved conditional laws. No theorem here
asserts that real IQ data satisfy the assumptions.\\
\end{longtable}
\endgroup

The shared probability-count results \ref{sm:binomial}, \ref{sm:wilson},
\ref{sm:clopper-pearson}, \ref{sm:conditional-count} and \ref{sm:rare}
are proved for completeness and use in the other two simulation reports.
This sequential report uses means and SDs of all generated outcomes,
rather than accepted-event proportions or random conditional-count
denominators. Its interval claims use the Gaussian sample results above.
